% Source: https://artofproblemsolving.com/community/c3h1450002_upvotes_geometry_questions
% Stable identity: aops-topic-1450002-q01
Joining the centre $O$ to the eight vertices cuts the octagon into $8$ congruent isosceles triangles with two sides $16$ and apex angle $\dfrac{360^\circ}{8}=45^\circ$. Each has area
\[
\dfrac{1}{2}\cdot16\cdot16\cdot\sin45^\circ=128\cdot\dfrac{\sqrt2}{2}=64\sqrt2 .
\]
So the octagon has area $8\cdot64\sqrt2=\boxed{512\sqrt2}$ square inches $\approx724.1$ square inches.
\begin{asy}
usepackage("amsmath");
size(5cm);
defaultpen(fontsize(10pt));
real rr=3; pair O=(0,0);
draw(circle(O,rr),gray);
pair[] V; for(int k=0;k<8;++k) V.push(rr*dir(22.5+45*k));
fill(O--V[0]--V[1]--cycle,palegreen);
draw(V[0]--V[1]--V[2]--V[3]--V[4]--V[5]--V[6]--V[7]--cycle,linewidth(0.9));
draw(O--V[0],blue); draw(O--V[1],blue);
draw(arc(O,0.6,22.5,67.5),red);
dot(O); for(int k=0;k<8;++k) dot(V[k]);
label("$O$",O,S); label("$45^\circ$",O+1.2*dir(45),red);
label("$16$",O+0.5*rr*dir(22.5),SE,blue);
label("\mbox{}",(min(currentpicture,user=true)+max(currentpicture,user=true))/2,black); // keeps the text after the figure black when LaTeX typesets the labels inline
\end{asy}

(Without trigonometry: the altitude from one vertex of the shaded triangle to the opposite radius makes a $45^\circ$-$45^\circ$-$90^\circ$ triangle with hypotenuse $16$, so this altitude is $8\sqrt2$ and the area is $\dfrac{1}{2}\cdot16\cdot8\sqrt2=64\sqrt2$.)
